\documentclass[../../../include/open-logic-section]{subfiles}

\begin{document}
\olfileid{sth}{ord-arithmetic}{intro}

\olsection{Pambuka}

Ing \olref[ordinals][]{chap}, kita wis ngembangaké téori wilangan ordinal.
Ing \olref[spine][]{chap}, kita weruh yèn ordinal bisa dianggep minangka
balung geger kang dadi sumbu pambangunan pérangan liya saka hierarki.
Nanging, iku dudu siji-sijiné peran ordinal. Ana uga tugas nindakake
aritmetika ordinal.

Bab iki wis tau kita singgung ing
\olref[ordinals][idea]{sec}, nalika ngrembug $\omega$,
$\omega+1$ lan $\omega+\omega$. Nalika iku, rembugan kita isih informal;
saiki wektuné njlentrehaké kanthi trep. Nanging, perlu disebutaké
yèn bab iki ora akèh ngrembug filsafat; isiné pangembangan teknis
lan sawijining pengamatan kang rada narik kawigaten: perkara kang
padha bisa ditindakake kanthi rong cara béda.

\end{document}
